\section*{Oppgave 3 - Layout}
\setcounter{section}{3}
\addcontentsline{toc}{section}{Oppgave 3 — Layout}

\setcounter{subsection}{0}

\subsection{Kviles layout}

Layout beskriver hvordan produksjonsressurser er fysisk plassert i fabrikken \parencite{strandhagen2024}. Hovedskillet går mellom \textit{funksjonell layout} (gruppert etter prosesstype) og \textit{linje-/produktlayout} (fast rekkefølge) \parencite{slack2019}.

Produksjonen hos Kvile består av aluminiumbearbeiding, snekkeravdeling, trekomponentmontasje, lakkering (to separate anlegg med betydelig omstillingstid), tørking, endemontering og pakking. Ulike produkter følger ulike ruter: rene treprodukter hopper over aluminiumbearbeiding, mens aluminium/tre-modellene behandles i parallell gjennom begge materialløpene før de møtes i endemonteringen. Dette samsvarer med \textit{funksjonell layout}: ressurser gruppert etter prosesstype, og produktene flyttes mellom avdelingene etter behov.

Lakkeanleggene har imidlertid trekk av \textit{linjelayout} internt, da alle deler følger samme sekvens (forbehandling $\rightarrow$ lakkering $\rightarrow$ tørking). Kviles layout kan dermed beskrives som en \textbf{kombinasjon}: primært funksjonell layout på fabrikknivå, med linjeorganiserte lakkeanlegg som delte ressurser. Figur~\ref{fig:layout-flyt} illustrerer materialflyten gjennom fabrikken.

Dersom Kvile går over til ATO for butikkmarkedet, som foreslått i oppgave~2, bør layouten også tilpasses. Komponentlageret bør plasseres nær endemontering og pakking, slik at sluttmontasjen kan gjennomføres med kort gjennomløpstid. Dette vil redusere intern transport og gjøre det enklere å styre montasjen etter faktiske kundeordrer.

\begin{figure}[htbp]
    \centering
    \begin{tikzpicture}[
        node distance=1.1cm,
        block/.style={rectangle, draw, rounded corners, minimum width=2.8cm, minimum height=0.85cm, align=center, font=\small},
        alu/.style={block, fill=blue!12},
        tre/.style={block, fill=green!12},
        felles/.style={block, fill=orange!15},
        arr/.style={-{Stealth[length=2.5mm]}, thick},
    ]
        % --- Venstre kolonne: Aluminium ---
        \node[alu] (alu-raw) {Aluminium (råvare)};
        \node[alu, below=of alu-raw] (alu-bearb) {Al-bearbeiding\\{\footnotesize kapp/bøy/slip/vask}};
        \node[alu, below=of alu-bearb] (alu-lakk) {Lakkering Al\\+ tørking (2 døgn)};
 
        % --- Høyre kolonne: Tre ---
        \node[tre, right=2.5cm of alu-raw] (tre-raw) {Treverk (råvare)};
        \node[tre, below=of tre-raw] (tre-bearb) {Snekkeravd.\\{\footnotesize kapp/fres/slip}};
        \node[tre, below=of tre-bearb] (tre-mont) {Trekomp.-montasje\\{\footnotesize plate/sete/rygg}};
        \node[tre, below=of tre-mont] (tre-lakk) {Lakkering tre\\+ tørking (2 døgn)};
 
        % --- Felles bunn ---
        \node[felles, below=1.3cm of $(alu-lakk)!0.5!(tre-lakk)$] (ende) {Endemontering};
        \node[felles, below=of ende] (pakk) {Pakking / lager};
 
        % Piler aluminium (ned)
        \draw[arr, blue!60] (alu-raw) -- (alu-bearb);
        \draw[arr, blue!60] (alu-bearb) -- (alu-lakk);
        \draw[arr, blue!60] (alu-lakk) |- (ende);
 
        % Piler tre (ned)
        \draw[arr, green!50!black] (tre-raw) -- (tre-bearb);
        \draw[arr, green!50!black] (tre-bearb) -- (tre-mont);
        \draw[arr, green!50!black] (tre-mont) -- (tre-lakk);
        \draw[arr, green!50!black] (tre-lakk) |- (ende);
 
        % Endemontering til pakking
        \draw[arr] (ende) -- (pakk);
 
        % Kolonne-labels
        \node[above=0.3cm of alu-raw, font=\footnotesize\bfseries, blue!70] {Aluminiumløp};
        \node[above=0.3cm of tre-raw, font=\footnotesize\bfseries, green!50!black] {Treløp};
 
    \end{tikzpicture}
    \caption[Materialflyt hos Kvile AS]{Materialflyt hos Kvile AS. Blått løp viser aluminiumbearbeiding, grønt løp viser trebearbeiding. De to løpene møtes i endemonteringen. Layouten er primært funksjonell, med linjeorganiserte lakkeanlegg.}
    \label{fig:layout-flyt}
\end{figure}